% SPDX-License-Identifier: CC-BY-4.0
% Copyright (c) 2026 Martin Schröder <info@swedishembedded.com>

\section{The first registered campaign: results}\label{sec:results}

Every number in this section is produced from the raw records in \code{data/runs}
by \code{scripts/build_data.py}. The experiments are pilots. E1 ran as designed; E2 and E3
ran but their registered criterion was not met: the registered runs produced no release, so
what they measure is where the loop stopped. The results are reported as registered and
measured; later repairs (\cref{sec:corrections}) do not rewrite them.

\subsection{E1: the funnel on a real corpus}\label{sec:e1}

The generator (Qwen3-4B) was asked for four kinds of task from the three pages of
Brain's documentation. The task report lists \mEOneTasks{} admitted tasks, covering
\mEOneParts{} of the \mEOnePartsTotal{} text parts of the corpus,
produced in \mEOneElapsedSecs{} seconds.\cl{M01} By kind,
recall admitted \mEOneAdmittedRecall, explain \mEOneAdmittedExplain, classify
\mEOneAdmittedClassify{} and the code-built denoise \mEOneAdmittedDenoise.\cl{M02}
The generator's proposals were rejected \mEOneRejectedRecall, \mEOneRejectedExplain,
\mEOneRejectedClassify{} and \mEOneRejectedDenoise{} times in the same order.\cl{M03}
\Cref{tab:e1} gives the rejections by reason. Of \mEOneRejectedTotal{} rejections,
\mEOneReasonNoSubject{} carry one reason, \code{no\_subject}: the generator did not
name what the task is about, which the admission rule of \cref{sec:admission}
requires. One was a duplicate and one was ungrounded.\cl{M04}
This is the first wave of E1 (a binary built before the repairs of \cref{sec:defects}).

\paragraph{Why the subject is missing.} We read each of the 82 first-wave rejections, and
those of the second-wave reruns below, against the rule of \cref{sec:admission}
(\code{scripts/classify_rejections.py}). The rule asks that the subject the generator names
be in the instruction and be named by the source or a cited section. Four classes result:
the subject is absent from the instruction (for example subject \code{braintop}, question
``What does the \code{--cli} option print?''); its words are in the instruction but not
adjacent; it is in the instruction but neither the source nor the cited section carries it
(an invented label such as ``metric line''); or it is in the instruction once code quotes
are ignored. The last class (18 of the 82, subject ``q key'' against ``the \code{q} key'')
was a defect of the rule, a literal comparison that markup defeats, and is repaired.

\paragraph{Second wave.} We reran E1 once per corpus on the repaired binary
(\cref{tab:e1w2}); wall times are not comparable between runs because the card was shared.
On the same three Brain pages the generator made \mEOneaProposals{} proposals, \mEOneaAdmittedProposals{} were
admitted and \mEOneaReasonNoSubject{} rejected as \code{no\_subject}
(\mEOneaShareNoSubject), of which \mEOneaClassAbsent{} name a subject absent from the instruction
and none is markup only.\cl{M18,M19} On three pages of the Sven user documentation it made
\mEOnebProposals{} proposals, admitted \mEOnebAdmittedProposals{} and rejected \mEOnebReasonNoSubject{}
as \code{no\_subject} (\mEOnebShareNoSubject), again with no markup-only case.\cl{M20}
The repair therefore removed the one class that was the rule's fault and did not change how
often the generator's subject is refused. With Qwen3-8B as the generator on the Brain pages the
share was \mEOnecShareNoSubject{} (\mEOnecProposals{} proposals), the interval overlapping that of
the 4B generator, and \mEOnecClassAbsent{} of the \mEOnecReasonNoSubject{} again named a subject
absent from the instruction.\cl{M21} We conclude only that the dominant cause is the shape of
the generator's output, a question that omits the subject it names, at both sizes; we did not
test a changed brief, so the prompt is not excluded as a cause.

\begin{table}[t]
\centering
\small
\input{generated/e1w2_table}
\caption{E1, second wave: proposals and their fate by run. The last column splits the
\code{no\_subject} rejections into subject absent from the instruction, its words apart, not
in the source, and markup only.}\label{tab:e1w2}
\end{table}

\begin{table}[t]
\centering
\input{generated/e1_table}
\caption{E1: admissions and rejections of the task funnel (Qwen3-4B generator, three
documentation pages, one run).}\label{tab:e1}
\end{table}

\subsection{E2 and E3: a turn of the loop on invented facts}\label{sec:e2}

\paragraph{First wave.} The first wave taught 24 facts per generation from one file per
generation, with the policy Qwen3-0.6B. The unadapted policy answered \mETwoBaseGZero{} of 24 facts
of generation~0, \mETwoBaseGOne{} of generation~1 and \mETwoBaseGTwo{} of generation~2 (at least
three of four probes right) and \mETwoBaseAnchor{} of 40 anchor questions.\cl{M05}
The first \code{learn} reached training with \mETwoGZeroTasks{} tasks admitted, none of the
\mETwoGZeroAttempts{} closed-book attempts verified correct, the teacher's answer verified for
\mETwoGZeroTeacherVerified{} tasks, and \mETwoGZeroDatasetRecords{} records; the second ended with
\mETwoGOneDatasetRecords{} records.\cl{M06,M07} Both exited with status \mETwoExitLearnGZero{} at the
training stage with ``\mETwoGZeroRefusal''.\cl{M08} The live policy was unchanged, no release was made
(\mETwoReleasesFinal{} listed) and the rollback was refused: ``\mETwoRollbackMessage''.\cl{M09,M10}
The cause is the first defect of \cref{sec:defects}: the single fact sheet was one group of source
text, which the holdout rule cannot divide. The rollback refusal was correct, since no release
existed. This wave measured nothing about learning, retention or release.

\paragraph{Second wave, registered run.} The facts are 30 per generation, one sheet file each,
with the training flags fixed from two earlier pilots on other seeds (\code{prereg/protocol.md}).
The unadapted policy answered \mETworBaseGZero{}, \mETworBaseGOne{} and \mETworBaseGTwo{} of the
30 facts of the three generations and \mETworBaseAnchor{} of 40 anchor questions, the chance level for
numeric answers being zero facts at three of four probes.\cl{M22}
The first \code{learn} admitted \mETworGZeroTasks{} tasks, of which \mETworGZeroTaught{} were taught
(the policy never solved them closed-book and the teacher's answer verified) and \mETworGZeroNever{}
were not (the teacher's answer failed verification); it wrote \mETworGZeroVariants{} variants,
\mETworGZeroRecords{} records, \mETworGZeroTrainRecords{} of them to train on after the held-out
families were taken out, and trained \mETworGZeroSteps{} steps at rank \mETworGZeroRank.\cl{M23}
Held-out whole families are facts the policy is never taught; that is the price of a
measurement independent of training (\cref{sec:defects}).

\Cref{tab:gate} gives the gate's four checks. The generation~0 candidate passed three: improvement
(it won \mETworGZeroImpSignWins{} of \mETworGZeroImpSignDiscordant{} discordant families, sign test
$p=\mETworGZeroImpSignP$, and \mETworGZeroImpVariantsRight{} of \mETworGZeroImpVariantsN{} paraphrase
tasks right against none for the champion), retention (nothing earlier to hold) and serve
(\mETworGZeroServeAgreed{} of \mETworGZeroServeSampled{} sampled answers of the served adapter
agreed with the in-process ones). It was refused by the anchor check: the 40-question anchor suite
fell by \mETworGZeroAnchorDrop{} against a bound of \mETworGZeroAnchorBound, which is one question
lost.\cl{M24} The release was therefore not made. The registered success criterion (a release and a McNemar $p<0.05$ against the base)\cl{RC04} was therefore
\emph{not met}: the outcome is ``no release'', for the reason given.

Measured on the probes (\cref{tab:e2}), the refused candidate answered \mETworCandidateGZeroGZero{} of
the 30 facts of generation~0 against none at base (\mETworBaseCandidateGZeroGZeroGained{} gained,
\mETworBaseCandidateGZeroGZeroLost{} lost, exact McNemar $p=\mETworBaseCandidateGZeroGZeroP$), answered
\mETworCandidateGZeroGOne{} of generation~1 and \mETworCandidateGZeroGTwo{} of the never-taught generation~2,
and \mETworCandidateGZeroAnchor{} of 40 anchor questions.\cl{M25}
The two lost anchor answers were a number written where a general fact was asked (``There are 17
continents'') and a wrong ocean. Learning the invented facts moved the policy where nothing taught
it to move, and the gate caught it; no other measurement here sees that.

The second generation could not continue a champion, since none existed: its \code{learn} trained from the
base on \mETworGOneTasks{} tasks (\mETworGOneTaught{} taught, \mETworGOneTrainRecords{} records trained on)
and met the same pattern: improvement, retention and serve passed (\mETworGOneServeAgreed{} of
\mETworGOneServeSampled{} served answers agreed), the anchor fell by \mETworGOneAnchorDrop{}, no release.\cl{M26}
Its candidate answered \mETworCandidateGOneGOne{} of the generation~1 facts (\mETworBaseCandidateGOneGOneGained{}
gained, \mETworBaseCandidateGOneGOneLost{} lost, $p=\mETworBaseCandidateGOneGOneP$) and \mETworCandidateGOneAnchor{}
of 40 anchor questions.\cl{M27} Because no first release existed, retention across two releases and rollback
to a first release are \emph{not measured}: the live policy stayed the base after each command
(\mETworAfterGZeroGZero{} facts and \mETworAfterGOneAnchor{} anchor answers throughout, \mETworReleasesFinal{} releases
listed) and the rollback was refused with status \mETworRollbackExit: ``\mETworRollbackMessage''.\cl{M28}

\begin{table}[t]
\centering
\small
\input{generated/gate_table}
\caption{E2 and E3, second wave: the release gate's four checks on each candidate. A refused
check is the reason the candidate was not released.}\label{tab:gate}
\end{table}

\begin{table}[t]
\centering
\small
\input{generated/e2w2_table}
\caption{E2 and E3, second wave: facts answered (at least three of four probes) with Wilson 95\%
intervals, and anchor questions right, for the unadapted policy and for each candidate
(measured whether or not the gate released it). 30 facts per generation.}\label{tab:e2}
\end{table}

\paragraph{E2b (post hoc).} Since the anchor check alone refused the candidate, we tried one other setting
decided after seeing that result: rehearsal of the base's own answers at 0.5 of the draws and a lower
rate, on generation~0 only. The candidate passed improvement (sign test $p=\mETwobGZeroImpSignP$), retention
and serve (\mETwobGZeroServeAgreed{} of \mETwobGZeroServeSampled{} agreed) and was refused by the anchor check
again, which fell by \mETwobGZeroAnchorDrop{}; it answered \mETwobCandidateGZeroGZero{} of the 30
generation~0 facts (\mETwobBaseCandidateGZeroGZeroGained{} gained, none lost,
$p=\mETwobBaseCandidateGZeroGZeroP$) and \mETwobCandidateGZeroAnchor{} of 40 anchor questions.\cl{M29}
More rehearsal did not protect the anchor here.

\paragraph{Repaired-verifier rerun.} Reading the teacher's failed verdicts showed that the stated-reference
verifier failed correct answers written with a thousands separator (``4,711'' for 4711), which, with
references that carry a unit word the answer omits, left \mETworGZeroNever{} of the registered run's
\mETworGZeroTasks{} tasks untaught (\cref{sec:defects}). We repaired it and, as
registered, reran the same design with a new seed. The rerun taught \mETwocGZeroTaught{} of \mETwocGZeroTasks{} tasks
(\mETwocGZeroNever{} not) and trained on \mETwocGZeroTrainRecords{} records. Improvement (sign test
$p=\mETwocGZeroImpSignP$; \mETwocGZeroImpVariantsRight{} of \mETwocGZeroImpVariantsN{} paraphrase tasks right) and
retention passed; the anchor check refused it (a drop of \mETwocGZeroAnchorDrop); and the serve check is
\emph{not measured} in this run, because the shell it was started from lacked the CUDA runtime path
and the served model faulted (the serve check passed in the other runs).\cl{M31} The candidate answered \mETwocCandidateGZeroGZero{}
of 30 facts of generation~0 against none at base (\mETwocBaseCandidateGZeroGZeroGained{} gained, none lost,
$p=\mETwocBaseCandidateGZeroGZeroP$), none of generations 1 and 2, and \mETwocCandidateGZeroAnchor{} of 40 anchor
questions; the generation~1 candidate, trained from the base because no champion existed, answered
\mETwocCandidateGOneGOne{} of the generation~1 facts and \mETwocCandidateGOneAnchor{} anchor questions, and was refused
with an anchor drop of \mETwocGOneAnchorDrop. The rollback was refused for lack of a release (status
\mETwocRollbackExit; \mETwocReleasesFinal{} releases listed).\cl{M32}

\paragraph{E2d (post hoc): a 1.7B policy.} The anchor check alone had refused every 0.6B candidate, so after seeing that
we asked, outside the registered design, whether a larger policy keeps the anchor while learning the same facts: Qwen3-1.7B,
the registered flags, a new seed, one run. Its first \code{learn} admitted \mETwodGZeroTasks{} tasks, taught \mETwodGZeroTaught{}, trained
on \mETwodGZeroTrainRecords{} records, and passed all four checks: improvement (\mETwodGZeroImpSignWins{} of \mETwodGZeroImpSignDiscordant{}
discordant families won, $p=\mETwodGZeroImpSignP$; \mETwodGZeroImpVariantsRight{} of \mETwodGZeroImpVariantsN{} paraphrase tasks right),
retention, anchor (a drop of \mETwodGZeroAnchorDrop{} against \mETwodGZeroAnchorBound) and serve (\mETwodGZeroServeAgreed{} of
\mETwodGZeroServeSampled{} answers agreed); the candidate was released.\cl{M33} Asked through the alias, the released policy answered
\mETwodAfterGZeroGZero{} of the 30 generation~0 facts against none at base (\mETwodBaseCandidateGZeroGZeroGained{} gained, none lost,
$p=\mETwodBaseCandidateGZeroGZeroP$), none of the facts of generations 1 and 2, and \mETwodAfterGZeroAnchor{} of 40 anchor
questions.\cl{M34} The second generation continued this champion with replay. Its candidate answered \mETwodCandidateGOneGOne{} of the
generation~1 facts but only \mETwodCandidateGOneGZero{} of the 15 generation~0 facts the champion answered, and \mETwodCandidateGOneAnchor{} of
40 anchor questions. The gate's retention check passed, on \mETwodGOneRetentionTasks{} tasks of which \mETwodGOneRetentionRight{} were right for
either model, because it holds the candidate to the earlier release's \emph{held-out} tasks, which for a facts campaign are facts that release
never taught; the anchor check refused the candidate (a drop of \mETwodGOneAnchorDrop), so the alias stayed on the first release
(\mETwodAfterGOneGZero{} facts, \mETwodAfterGOneAnchor{} anchor answers).\cl{M35} The rollback was refused, as specified, because the first release has
no predecessor.\cl{M36} Retention across two releases was therefore measured only as a refused candidate's forgetting, on independent probes,
not as a release; the gate's own retention check did not see that forgetting. The first release's probes through \code{brain serve}
beyond the gate's sample of eight were not measured: the server refused requests under load on the shared card.

\paragraph{Re-grading.} After the runs the stated-reference verifier was repaired so that a number
written with a thousands separator is the same number (\cref{sec:defects}). The candidates of the
registered run and of E2b were measured again with the repaired verifier and gave the same counts
(\mETworRegradeCandidateGZeroGZero{} and \mETworRegradeCandidateGOneGOne{} facts, and
\mETwobRegradeCandidateGZeroGZero{}; anchor \mETworRegradeCandidateGZeroAnchor, \mETworRegradeCandidateGOneAnchor{} and
\mETwobRegradeCandidateGZeroAnchor).\cl{M30} The verifier's effect was on training, not on these probes:
it decided which facts the teacher's answers could teach.
